\documentclass[11pt,a4paper]{article}
\usepackage[top=1.7cm,bottom=1.5cm,left=1.5cm,right=1.5cm,headsep=8pt,footskip=16pt]{geometry}
\usepackage{enumitem}
\usepackage{amssymb}
\usepackage{array}
\usepackage[table]{xcolor}
\usepackage{fancyhdr}
\usepackage{lastpage}
\usepackage{multicol}
\usepackage{polyglossia}
\setdefaultlanguage[calendar=gregorian, numerals=maghrib]{arabic}
\setotherlanguage{english}
\IfFontExistsTF{Amiri}{%
  \newfontfamily\arabicfont[Script=Arabic, Scale=1.08]{Amiri}
}{%
  \newfontfamily\arabicfont[Script=Arabic, Scale=1.08, Path=../fonts/, Extension=.ttf,
    UprightFont=*-Regular, BoldFont=*-Bold, ItalicFont=*-Italic,
    BoldItalicFont=*-BoldItalic]{Amiri}
}
\IfFontExistsTF{AlBattar}{%
  \newfontfamily\khatt[Script=Arabic, Scale=1.15]{AlBattar}
}{%
  \newfontfamily\khatt[Script=Arabic, Scale=1.15, Path=../fonts/,
    BoldFont=ae_AlBattar.ttf]{ae_AlBattar.ttf}
}
\renewcommand{\headrulewidth}{0.8pt}
\renewcommand{\footrulewidth}{0.8pt}
\fancyhf{}
\rhead{\khatt المدرسة العليا للأساتذة ببوزريعة}
\lhead{\khatt ملحقة بومرداس}
\lfoot{\khatt الأستاذ: د. بن عبد الرزاق يوسف}
\cfoot{\thepage/\pageref{LastPage}}
\pagestyle{fancy}
\parindent 0pt
\setlength{\columnsep}{0.6cm}
\setlength{\columnseprule}{0.4pt}
\setlist[enumerate,1]{label=\textbf{\arabic*.}, leftmargin=*, itemsep=5pt, topsep=2pt}
\begin{document}

\linespread{1.08}\selectfont
\begin{center}
{\LARGE\khatt\textbf{الاختبار الرابع عشر في الذكاء الاصطناعي}}\\[4pt]
{\normalsize وضعيات إشكالية ــ أسئلة متعددة الاختيارات ــ الفصول من 1 إلى 6}
\end{center}
\vspace{-2pt}\rule{\textwidth}{0.6pt}\\[2pt]
\textbf{التعليمة:} يتضمن الاختبار خمس وضعيات إشكالية. اقرأ كل وضعية بعناية، ثم اختر التحليل أو التصرف الأنسب
من بين الاختيارات الخمسة؛ لكل وضعية إجابة صحيحة واحدة فقط.
اكتب \textbf{حرف} الإجابة الصحيحة في جدول الإجابات في آخر الاختبار.
لكل وضعية 4 نقاط (المجموع 20 نقطة).
\vspace{4pt}

\noindent
{\large\khatt\textbf{الوضعية 1}}\hfill{\small(الفصل 1)}\par\nopagebreak\vspace{3pt}
\setlength{\fboxsep}{7pt}\colorbox{gray!10}{\parbox{\dimexpr\linewidth-14pt}{كتب تلميذ قصة بمساعدة أداة توليدية، ثم سأل أستاذته: «هل الأداة تفهم القصة كما أفهمها أنا؟ لقد اقترحت نهاية مؤثرة جدًا». وتريد الأستاذة جوابًا دقيقًا يناسب سنّه، ويبيّن حدود الأداة دون تهويل أو تهوين. وتتوقع أن يطرح تلاميذ آخرون السؤال نفسه في الحصص القادمة، فتريد جوابًا موحدًا. أي جواب هو الأدق؟}}\par\nopagebreak
\begin{itemize}[leftmargin=1.0cm, labelsep=6pt, itemsep=3pt, topsep=4pt, parsep=0pt]
  \item[\textbf{أ)}] الأداة اختارت النهاية عشوائيًا تمامًا دون أي صلة بما تعلّمته، ولذلك فإن جودة اقتراحها مصادفة لا تتكرر، ولا فائدة من الاستعانة بها مرة أخرى في الكتابة.
  \item[\textbf{ب)}] الأداة تعلّمت من نصوص كثيرة كيف تُكتب النهايات المؤثرة عادة، فتنتج ما يشبهها، لكنها لا تعيش المشاعر ولا التجارب كما يعيشها الإنسان؛ فاقتراحها مفيد، والحكم على القصة لك.
  \item[\textbf{ج)}] لا، الأداة لا تفعل شيئًا سوى نسخ نهايات قصص موجودة حرفيًا من الإنترنت، فالنهاية المقترحة منقولة بالضرورة من قصة أخرى، ولا يجوز لك استعمالها في قصتك، وعليك أن تبحث عن نهاية أخرى بنفسك.
  \item[\textbf{د)}] الأداة تفهم القصة أفضل منك، لأنها قرأت ملايين القصص، ولذلك يجب أن تقبل كل ما تقترحه دون تعديل، فهي أقدر منك على الحكم على ما يؤثر في القرّاء.
  \item[\textbf{هـ)}] نعم، الأداة تفهم القصة وتشعر بها مثلك تمامًا، ولذلك اقترحت نهاية مؤثرة، فهي تملك مشاعر حقيقية تعلّمتها من الكتب، ويمكنك أن تثق بحكمها على القصة أكثر من حكمك.
\end{itemize}
\par\vspace{6pt}
\noindent
{\large\khatt\textbf{الوضعية 2}}\hfill{\small(الفصل 2)}\par\nopagebreak\vspace{3pt}
\setlength{\fboxsep}{7pt}\colorbox{gray!10}{\parbox{\dimexpr\linewidth-14pt}{لاحظ أستاذ أن أداة ذكاء اصطناعي تكتب نصوصًا عربية سليمة، لكن بعضها يبدو «مترجمًا»: تراكيب مثل «يلعب دورًا مهمًا في»، و«في نهاية اليوم» بمعنى «في المحصلة». وفسّر كل زميل الظاهرة بطريقته. ويريد أن يقدّم لطلبته إرشادات لمراجعة أسلوب النصوص المولَّدة. أي تفسير هو الأقرب إلى طريقة تعلّم النماذج؟}}\par\nopagebreak
\begin{itemize}[leftmargin=1.0cm, labelsep=6pt, itemsep=3pt, topsep=4pt, parsep=0pt]
  \item[\textbf{أ)}] هذه التراكيب هي الفصحى الصحيحة، والأسلوب العربي القديم هو الذي تجاوزه الزمن، فالنموذج تعلّم اللغة المعاصرة الصحيحة، ولا حاجة إلى أي مراجعة أسلوبية.
  \item[\textbf{ب)}] النموذج يفكّر بالإنجليزية ثم يترجم كل جملة حرفيًا قبل عرضها، ولذلك يستحيل أن يكتب بأسلوب عربي أصيل مهما تدرّب على نصوص عربية أصيلة.
  \item[\textbf{ج)}] النموذج تعلّم من نصوص عربية أصيلة فقط، لكنه يُدخل هذه التراكيب عمدًا ليبدو حديثًا، لأن النماذج مصممة لتقليد أسلوب وسائل الإعلام المعاصرة في كل النصوص، وهذا ما يجعل نصوصها قريبة من لغة الصحافة.
  \item[\textbf{د)}] الظاهرة سببها أن الأستاذ كتب طلبه بأسلوب مترجم، فالنموذج لا يتأثر ببيانات تدريبه بل بأسلوب الطلب وحده، ولو كتب طلبًا فصيحًا لجاء النص خاليًا منها تمامًا.
  \item[\textbf{هـ)}] النصوص الإنجليزية تغلب على بيانات تدريب كثير من النماذج، والنصوص العربية المتاحة فيها كثير من الترجمات، فتتعلّم النماذج هذه التراكيب وتعيد إنتاجها؛ ويعالجها المستخدم بالمراجعة الأسلوبية.
\end{itemize}
\par\vspace{6pt}
\noindent
{\large\khatt\textbf{الوضعية 3}}\hfill{\small(الفصل 4)}\par\nopagebreak\vspace{3pt}
\setlength{\fboxsep}{7pt}\colorbox{gray!10}{\parbox{\dimexpr\linewidth-14pt}{تستعمل أستاذة أداة ذكاء اصطناعي لمساعدتها في تقدير علامات الإنشاءات، فلاحظت أن الأداة تمنح باستمرار علامات أعلى للنصوص الأطول والمليئة بالمفردات النادرة، حتى إن كانت أفكارها مفككة. والعلامات ستُرصد في كشوف رسمية. وعدد الإنشاءات يتجاوز مئة وخمسين في كل فصل، ووقتها محدود. ما التصرف الأنسب؟}}\par\nopagebreak
\begin{itemize}[leftmargin=1.0cm, labelsep=6pt, itemsep=3pt, topsep=4pt, parsep=0pt]
  \item[\textbf{أ)}] تخصم آليًا نقطتين من كل نص يتجاوز صفحة واحدة، لتعويض انحياز الأداة نحو الطول، ثم تعتمد العلامات الناتجة دون أي مراجعة فردية للنصوص.
  \item[\textbf{ب)}] تعتمد علامات الأداة، لأن الطول والمفردات النادرة دليل على الجهد والثراء اللغوي، وتفكك الأفكار عيب ثانوي لا يؤثر في جودة الإنشاء عند تقويمه.
  \item[\textbf{ج)}] تطلب من التلاميذ كتابة نصوص أطول بمفردات نادرة، لأن ذلك يرفع علاماتهم عند الأداة، فتتحقق العدالة بتكيّف التلاميذ مع معايير الأداة لا العكس.
  \item[\textbf{د)}] تتنبّه إلى أن الأداة منحازة إلى الطول والمفردات النادرة، فتزوّدها بشبكة تقويم صريحة المعايير، وتراجع علاماتها بنفسها، خاصة للنصوص القصيرة المحكمة، لأن العلامة النهائية مسؤوليتها.
  \item[\textbf{هـ)}] تزوّد الأداة بشبكة تقويم، ثم تعتمد علاماتها دون مراجعة، لأن الشبكة وحدها تزيل كل انحياز، ولا يمكن أن تعود الأداة إلى تفضيل الطول بعد تزويدها بالمعايير.
\end{itemize}
\par\vspace{6pt}
\noindent
{\large\khatt\textbf{الوضعية 4}}\hfill{\small(الفصل 5)}\par\nopagebreak\vspace{3pt}
\setlength{\fboxsep}{7pt}\colorbox{gray!10}{\parbox{\dimexpr\linewidth-14pt}{قرأ أستاذ مقالًا إلكترونيًا مقنعًا عن «فوائد تعليم النحو بالألعاب الرقمية»، كاتبه مجهول، والموقع يبيع تطبيقًا للألعاب التعليمية. ثم طلب من أداة ذكاء اصطناعي تقييم المقال، فوصفته بأنه «موثوق ومدعوم بالأدلة». ويريد أن يستعمل المقال في تكوين زملائه حول الألعاب التعليمية. ما المنهج الأنسب لتقييم موثوقية المقال؟}}\par\nopagebreak
\begin{itemize}[leftmargin=1.0cm, labelsep=6pt, itemsep=3pt, topsep=4pt, parsep=0pt]
  \item[\textbf{أ)}] يقيس موثوقية المقال بعدد مشاركاته في مواقع التواصل، لأن المقالات الأكثر تداولًا خضعت لتقييم جمهور واسع، فيكون الانتشار مؤشرًا كافيًا على صحة أدلتها.
  \item[\textbf{ب)}] يعتمد حكم الأداة، لأنها حلّلت المقال تحليلًا محايدًا، وبيع الموقع للتطبيق لا يؤثر في مصداقية المقال، فالمعلومة تُقيَّم بمضمونها وحده لا بمصدرها.
  \item[\textbf{ج)}] يقرأ قراءة جانبية: يبحث عن الكاتب والموقع ومصالحه التجارية، ويتتبع الأدلة المذكورة إلى دراسات منشورة، ويقارنها بمصادر مستقلة، لأن حكم الأداة على الموثوقية لا يغني عن فحص المصدر.
  \item[\textbf{د)}] يرفض المقال فورًا لأن الموقع تجاري، دون قراءة أدلته أو تتبعها، لأن كل مقال في موقع تجاري خاطئ بالضرورة، ولا يمكن أن يتضمن أي معلومة صحيحة، ويوفر عليه الوقت في تكوين الزملاء.
  \item[\textbf{هـ)}] يطلب من الأداة نفسها البحث عن الكاتب، فإن قالت إنه خبير معروف اعتمد المقال، لأن الأداة تعرف كل الكتّاب، ولا حاجة إلى التحقق من الدراسات المذكورة، ثم ينتقل مباشرة إلى استعمال المقال في التكوين.
\end{itemize}
\par\vspace{6pt}
\noindent
{\large\khatt\textbf{الوضعية 5}}\hfill{\small(الفصل 6)}\par\nopagebreak\vspace{3pt}
\setlength{\fboxsep}{7pt}\colorbox{gray!10}{\parbox{\dimexpr\linewidth-14pt}{أنتجت أستاذة بأداة توليد صوتي «حوارًا إذاعيًا» بين شخصيتين عن آداب الحوار، لتدريب التلاميذ على الاستماع. جاء الصوتان طبيعيين جدًا، لكن تلميذًا سأل: «هل هذان شخصان حقيقيان؟»، ولاحظت الأستاذة أخطاء في نطق بعض الكلمات. وتريد إعادة استعمال الحوار مع أقسام أخرى طوال السنة. ما التصرف الأنسب؟}}\par\nopagebreak
\begin{itemize}[leftmargin=1.0cm, labelsep=6pt, itemsep=3pt, topsep=4pt, parsep=0pt]
  \item[\textbf{أ)}] تخبر التلاميذ بوضوح أن الصوتين مولَّدان، وتصحح الأخطاء بإعادة توليد المقاطع أو تسجيلها بنفسها، وتستثمر السؤال في توعيتهم بأن الصوت الطبيعي لا يدل على وجود متحدث حقيقي.
  \item[\textbf{ب)}] تتخلى عن الحوار المولَّد نهائيًا، لأن طبيعية الأصوات المولَّدة خطر لا يمكن معالجته تربويًا، وأي استعمال لها في القسم يخدع التلاميذ مهما وقع التصريح بطبيعتها.
  \item[\textbf{ج)}] تخبر التلاميذ أن الصوتين مولَّدان، وتُبقي أخطاء النطق دون تصحيح، لأن الإفصاح عن طبيعة الصوت يغني عن مراجعة جودة النطق في أي نشاط تعليمي للاستماع.
  \item[\textbf{د)}] تتجاهل سؤال التلميذ وتصحح أخطاء النطق فقط، لأن طبيعة الصوت مسألة تقنية لا علاقة لها بأهداف الدرس، ولا يحتاج التلاميذ إلى معرفة مصدر الأصوات التي يستمعون إليها، فالمهم أن يتدرب التلاميذ على الاستماع.
  \item[\textbf{هـ)}] تؤكد للتلاميذ أنهما شخصان حقيقيان، لأن ذلك يزيد اهتمامهم بالنشاط، وتترك أخطاء النطق لأنها تشبه أخطاء المتحدثين الحقيقيين، فتجعل الحوار أكثر واقعية.
\end{itemize}
\par\vspace{6pt}
\begin{center}
{\khatt\textbf{جدول الإجابات}}\\[3pt]
\renewcommand{\arraystretch}{1.6}
\begin{tabular}{|>{\centering\arraybackslash}m{1.9cm}|*{5}{>{\centering\arraybackslash}m{1.2cm}|}}
\hline
\textbf{رقم السؤال} & \textbf{1} & \textbf{2} & \textbf{3} & \textbf{4} & \textbf{5} \\ \hline
\textbf{الإجابة} &  &  &  &  &  \\[8pt] \hline
\end{tabular}

\vspace{8pt}
\renewcommand{\arraystretch}{2.2}
\begin{tabular}{|>{\centering\arraybackslash}m{7.5cm}|>{\centering\arraybackslash}m{3.2cm}|>{\centering\arraybackslash}m{3.2cm}|}
\hline
\textbf{الاسم واللقب} & \textbf{الفوج} & \textbf{العلامة} \\ \hline
 & & \hfill /20 \\ \hline
\end{tabular}
\end{center}
\end{document}
